\Needspace{7\baselineskip}
\section{Count Flow and the Scoped C4 Invariant}
\label{sec:10}
Section 9 fixed the rule that presentation follows derivation. Section 10 applies that rule to the observer-side count. It derives the count first, records it separately from the core bookkeeping matrix, and then distinguishes it from a second occurrence of the same rational number in the intrinsic spectral construction.

\Needspace{5\baselineskip}
\subsection{Fixed Increment and Live Residual History}
The geometric residual share is the fixed exact quantity

\begin{equation*}
\begin{adjustbox}{max width=\linewidth,center}
$\displaystyle
\Omega
=
1-\frac{\pi}{3\sqrt2}.
$
\end{adjustbox}
\end{equation*}
A cell's live residual register is a different object. For a cell \(x\) with its own completed-update count \(K_x\),

\begin{equation*}
\begin{adjustbox}{max width=\linewidth,center}
$\displaystyle
S_{\mathrm{res},x}(K_x)
=
\operatorname{frac}
\left(
S_{\mathrm{res},x}(0)+K_x\Omega
\right).
$
\end{adjustbox}
\end{equation*}
Every cell obeys the same increment rule, but the accumulated states need not be equal. Their birth and completion histories may differ, and the retained residual record preserves that distinction. A common increment is therefore not a common live state.

\Needspace{5\baselineskip}
\subsection{The Four-Cell Reference Frame}
At the first four-cell closure, one designated cell occupies the observer/reference role. The other cells are registered relational content; they do not thereby become additional observers. The reference orients the reading of distinctions already present in the construction. It does not create the cells, their histories, or their closure relations.

The frame contains three cells read relative to the designated reference. Exactly one is co-present with it and content-indistinguishable from it at that stage; the remaining \(c-1\) cells are distinguished by the inherited completion order. Because the closure count is

\begin{equation*}
\begin{adjustbox}{max width=\linewidth,center}
$\displaystyle
c=3,
$
\end{adjustbox}
\end{equation*}
the content-distinguishable fraction is

\begin{equation*}
\begin{adjustbox}{max width=\linewidth,center}
$\displaystyle
S_{\mathrm{flow}}
:=
\frac{\text{content-distinguishable read cells}}
{\text{read cells}}
=
\frac{c-1}{c}.
$
\end{adjustbox}
\end{equation*}
% LAYOUT_ONLY_GUARD:STANDALONE_HENCE
\Needspace{7\baselineskip}
Hence

\begin{equation*}
\begin{adjustbox}{max width=\linewidth,center}
$\displaystyle
\boxed{
S_{\mathrm{flow}}
=
\frac{c-1}{c}
=
\frac23.
}
$
\end{adjustbox}
\end{equation*}
This is a count result: two distinguishable cells out of three read cells. It is derived rather than assigned. Choosing the other member of the co-present pair as the designated reference leaves the same one co-present cell and the same two content-distinguishable cells, so the count is independent of that permitted reference choice.

\Needspace{5\baselineskip}
\subsection{The Core Matrix and the Separate Invariant Record}
Section 4 already assembled

\begin{equation*}
\begin{adjustbox}{max width=\linewidth,center}
$\displaystyle
M_{\mathrm{cell}}
=
\begin{pmatrix}
I & c & \eta\\[2pt]
\tau & E & D\\[2pt]
m & G & \Omega
\end{pmatrix}.
$
\end{adjustbox}
\end{equation*}
Every entry is established independently before placement. This \(3\times3\) array remains the complete rectangular bookkeeping display used by this manuscript.

The observer-side count is recorded separately:

\begin{equation*}
\begin{adjustbox}{max width=\linewidth,center}
$\displaystyle
\boxed{
C_4^{\mathrm{inv}}
:=
S_{\mathrm{flow}}
=
\frac23.
}
$
\end{adjustbox}
\end{equation*}
This equation is not drawn as an attached fourth column. It introduces no lower service or support boxes and supplies no unfilled matrix offices. The manuscript consumes only this C4 invariant. It neither displays nor uses lower C4 carrier assignments; their broader interpretation and publication treatment lie outside the present derivation.

The notation \(C_4^{\mathrm{inv}}\) marks the provenance of the count. It does not turn the record into an algebraic matrix column, and no scientific inference follows from placing the equation before, after, above, or below the core matrix.

% LAYOUT_ONLY_HEADING_GUARD:SEPARATE_PROJECTOR_ACTION_VALUE
\Needspace{12\baselineskip}
\subsection{A Separate Projector/Action Value}
% LAYOUT_ONLY_GUARD:OBSERVER_PROJECTOR_SEQUENCE
\Needspace{20\baselineskip}
The intrinsic spectral construction contains another exact \(2/3\). Let

\begin{equation*}
\begin{adjustbox}{max width=\linewidth,center}
$\displaystyle
\Pi_o=E_0+E_1
$
\end{adjustbox}
\end{equation*}
% LAYOUT_ONLY_GUARD:OBSERVER_CYCLIC_AVERAGING
\Needspace{18\baselineskip}
be the rooted two-active projector on the three-state carrier. Cyclic averaging gives

\begin{equation*}
\begin{adjustbox}{max width=\linewidth,center}
$\displaystyle
\frac13
\sum_{r=0}^{2}
P^r\Pi_oP^{-r}
=
\frac23\mathbf1_V,
$
\end{adjustbox}
\end{equation*}
and the normalized support valuation gives

\begin{equation*}
\begin{adjustbox}{max width=\linewidth,center}
$\displaystyle
\nu(\Pi_o)=\frac23.
$
\end{adjustbox}
\end{equation*}
The associated one-pass action aperture and phase therefore satisfy

\begin{equation*}
\begin{adjustbox}{max width=\linewidth,center}
$\displaystyle
\boxed{
\nu(\Pi_o)
=
A_o^\#
=
\Phi_o
=
\frac23.
}
$
\end{adjustbox}
\end{equation*}
The flow count \(S_{\mathrm{flow}}\) and the projector/action value \(\nu(\Pi_o)\) have different inputs, types, and derivations. Their numerical equality does not identify the two typed objects. The charged-lepton spectral phase is supplied by the rooted projector/action construction, not by matrix placement, the C4 invariant, or the bookkeeping display.

Section 10 establishes the observer-side count, records it as a separate invariant, preserves the core \(3\times3\) matrix as an organizational display, and separates that count from the independently derived projector/action value. It does not extend the matrix, assign lower C4 roles, perform matrix algebra, or derive the spectral phase from the observer-side count.

